\documentstyle[%


righttag%


,11pt%


,amssymb%


,russian%               remove in Eng.


,izvru%                 Set izven% in Eng.


]{amsart}





% seek \text{...} !!


\newcommand{\const}{\operatorname{const}}


\newcommand{\mes}{\operatorname{mes}}


\newcommand{\tts}[1]{} 





\theoremstyle{plain}


\theorembodyfont{\it}


\newtheorem{thmnonum}{\hskip\parindent Теорема}


\renewcommand{\thethmnonum}{}








%\hoffset=-15mm%                  For Eng.


\setcounter{page}{45}


\god{2003}


\nomer{1~(488)} %                        Remove in Eng.


%\nomer{Vol.\,47, No.\,1}%               Set in Eng.


%\str{\pageref{begin}--\pageref{end}}%  Set in Eng.


\udk{517.958}





\author{\it О.А.\,Задворнов}


% NB:   ^^ remove in Eng.


\title{Постановка и исследование стационарной задачи\\ о контакте


мягкой оболочки с препятствием}


\thanks{Работа выполнена при финансовой поддержке 


Российского фонда фундаментальных исследований 


(коды проектов 01-01-00616, 03-01-00380).}





\date{}


%\pagestyle{myheadings}%         Set in Eng.


\pagestyle{plain} %              Remove in Eng.


\begin{document}


%\thispagestyle{plain}%          Set in Eng.


\maketitle


\label{begin}








\section*{Введение}





Рассматривается задача об определении положения


равновесия мягкой оболочки, закрепленной по краям, находящейся под 


воздействием массовой и поверхностной нагрузки и ограниченной в 


перемещении препятствием. Деформации и перемещения оболочки 


допускаются конечными.





Плоская задача о равновесном положении нити под воздействием нагрузки,  


ограниченной в расположении полуплоскостью, поставлена в виде вариационного


неравенства в \cite{BS_86}. Там же установлено существование решения этой


задачи. Случай, когда граница препятствия описывается вогнутой функцией, для


плоской задачи рассмотрен в \cite{BZB_2001_2}. В работе 


\cite{BBZ_2002_1} для этой же задачи


рассмотрен случай, когда множество допустимых конфигураций произвольно.





В данной работе рассматривается пространственная задача


о равновесном состоянии мягкой оболочки, при условии, что поверхность


препятствия описывается достаточно гладкой (не обязательно выпуклой)


функцией. Сначала, исходя из уравнений равновесия, записанных в 


декартовой системе координат, сформулирована поточечная задача. 


Затем на основе принципа виртуальных перемещений получена вариационная 


формулировка. Установлена эквивалентность указанных задач.





Во второй части статьи рассмотрен случай 


мягкой сетчатой оболочки, образованной


двумя семействами нитей. При определенных условиях на функции,


описывающие физические соотношения в нитях, поставлена обобщенная задача


в виде квазивариационного неравенства в банаховом пространстве и 


установлена ее разрешимость. Отметим, что соответствующая задача 


при отсутствии препятствия была изучена в \cite{BZ_92IV}.





\section{Поточечная и вариационная формулировки задачи}





Введем декартову систему координат $(x_1,x_2,x_3)$.


Считаем, что в недеформированном состоянии оболочка может быть описана 


поверхностью $\xi(\alpha)=(\xi_1(\alpha),\xi_2(\alpha),\xi_3(\alpha))$,


где $\alpha=(\alpha_1,\alpha_2) \in \Omega$ --- лагранжевы координаты,


$\Omega$ --- ограниченная область из $R^{2}$ с непрерывной 


по Липшицу границей $\Gamma$; предполагаем, что функция $\xi$


удовлетворяет условиям


\begin{equation}


\xi \in [C_1(\ov \Omega) ]^3,


\ \ | [\partial_1 \xi (\alpha),\partial_2 \xi (\alpha) ] |


\geq c > 0 \ \  \forall\alpha\in \ov \Omega.


\label{2}


\end{equation}





Через $w(\alpha)=(w_1(\alpha),w_2(\alpha),w_3(\alpha))$ обозначим функцию,


описывающую поверхность оболочки в деформированном состоянии;


$G(\alpha)={| [\partial_1 w(\alpha),\partial_2 w(\alpha) ] |}^2$ --- 


дискриминант метрического тензора поверхности 


деформированной оболочки.





Здесь использованы обозначения


$\partial_k={\partial}/{\partial \alpha_k}$, $k=1,2$;


$[\cdot,\cdot]$, $(\cdot,\cdot)$ и $| \cdot |$ --- 


векторное, скалярное произведения и норма в $R^{3}$ соответственно.





Уравнение равновесия оболочки, находящейся 


под воздействием внешних сил в отсутствии


препятствия, имеет в декартовой системе 


координат следующий вид \cite{R_Gul}:


\begin{equation}


\sum_{k,m=1}^{2} \partial_m (\sqrt G T^{km} \partial_k w )+


\sqrt G P+\sqrt G \gamma Q=0, 


\label{Ur_Rav}


\end{equation}


где $P$, $Q$ --- векторы плотности соответственно поверхностной и массовой 


нагрузок, $\gamma$ --- плотность материала оболочки в деформированном 


состоянии, $T^{km}$ --- ковариантные компоненты тензора напряжений


$$


{\mathrm T}=\sum_{k,m=1}^{2} T^{km}R_k R_m,


$$


$R_k(\alpha)=\partial_k w(\alpha)$, $k=1,2$, ---


векторы, образующие ковариантный 


локальный базис на деформированной поверхности. 





Далее будем считать, что расположение оболочки в пространстве 


ограничено препятствием. Взаимодействие препятствия с оболочкой


учтем, внеся в уравнение (\ref{Ur_Rav}) дополнительную поверхностную


нагрузку $P_0$ --- плотность силы реакции препятствия:


\begin{equation}


D(w)+\sqrt G P_0=0, 


\label{Ur_Rav_Pr_0}


\end{equation}


где 


$D(w)=\sum\limits_{k,m=1}^{2} \partial_m (\sqrt G T^{km} \partial_k w )+


\sqrt G P+\sqrt G \gamma Q$.


Такое моделирование взаимодействия тела с препятствием используется


при постановке контактных задач в теории упругости


(см. \cite{Duv}, \cite{Pan}).


Этот же подход использован в \cite{Kid} 


при исследовании стационарной осесимметричной задачи об эластичной 


цилиндрической оболочке, внутри которой содержится твердое тело.


Заметим, что изначально функция $P_0$ неизвестна и наряду с равновесным


положением оболочки $w$ является искомой.





Будем полагать в дальнейших рассуждениях, что входящие в 


уравнение (\ref{Ur_Rav_Pr_0}) функции обладают необходимой гладкостью.





Предположим, что поверхность препятствия задается в виде


$x_3=\varphi (x_1,x_2)$, где $\varphi$ ---


непрерывно-дифференцируемая функция,


и оболочка находится ``над препятствием'', т.\,е. \linebreak


$\xi_3(\alpha ) \geq \varphi (\xi_1(\alpha ),\xi_2(\alpha ))$,


$\alpha \in \ov\Omega$.


Для удобства дальнейшего изложения


введем функцию $F:R^3 \to R^1$ соотношением


$F(x_1, x_2,x_3)=\varphi (x_1,x_2)$.


Обозначим через $N:R^3 \to R^3$ вектор-функцию, связанную с единичной


внешней нормалью к поверхности препятствия по формуле 


\begin{equation} 


N(x)=\frac {\nu (x)} {|\nu (x)|},\ \ x \in R^3,


\label{Opr_N}


\end{equation}


где


$\nu=[\tau_1,\tau_2]$, 


$\tau_1=( 1,0,{\partial F}/{\partial x_1} )$,


$\tau_2=( 0,1,{\partial F}/{\partial x_2} )$.


Введем множество допустимых конфигураций оболочки


$M=\{v : \ov\Omega \to R^3$, $v_3(\alpha ) \geq F(v(\alpha ))$,


$\alpha \in \ov\Omega$, ${v} |_\Gamma={\xi} |_\Gamma\}$.


Функции из множества~$M$ описывают поверхности, находящиеся 


``над препятствием'' и удовлетворяющие граничному условию.





Будем считать материал препятствия абсолютно твердым,


а его поверхность --- абсолютно гладкой, т.\,е. препятствие при воздействии


на него не деформируется и порождает усилия только в направлении 


внешней нормали к своей поверхности.


Тогда плотность силы реакции препятствия можно представить в виде 


$P_0(\alpha )=\beta(\alpha ) N(w(\alpha ))$,


где $\beta : \ov\Omega \to R $ --- 


неизвестная функция, удовлетворяющая условиям


\begin{align} 


\beta(\alpha ) \geq &0,\ \ \alpha \in I(w); 


\label{beta_1} \\


%


\beta(\alpha )=&0,\ \  \alpha \in I^{-}(w).


\label{beta_2}


\end{align}


Здесь $I(w)=\{\alpha \in \ov\Omega : w_3(\alpha )=F(w(\alpha )) \}$ --- 


так называемое коинцидентное множество, 


$I^{-}(w)=\ov\Omega \setminus I(w)$. Заметим, что 


$I^{-}(w)=\{\alpha \in \ov\Omega : w_3(\alpha ) > F(w(\alpha )) \}$


--- множество тех точек, где оболочка


не контактирует с препятствием.





Таким образом, постановка задачи о равновесном положении закрепленной по краю


мягкой оболочки под воздействием нагрузки и ограниченной в пространстве


абсолютно твердым и гладким препятствием сводится к поиску функций


${w \in M}$ и $\beta$,


удовлетворяющих условиям (\ref{beta_1}), (\ref{beta_2})


и уравнению 


\begin{equation}


D(w(\alpha ))+\sqrt {G(\alpha )} \beta(\alpha ) N(w(\alpha ))=0,


\ \ \alpha \in \Omega.


\label{Ur_Rav_Pr}


\end{equation}





Перейдем от поточечной к вариационной формулировке этой задачи.


Для произвольного положения оболочки $u \in M$


определим множество допустимых направлений, достаточно малый 


сдвиг по которым из $u$ принадлежит $M$,


$$


M(u)=\{{v :\ov\Omega\to R^3,\ \exists t_v > 0;


\ \forall t \in [0,t_v],\ u+t v \in M} \} .


$$


Умножим скалярно уравнение


(\ref{Ur_Rav_Pr}) на произвольную функцию $\eta$ из множества $M(w)$ и


затем проинтегрируем по множеству $\Omega$


\begin{equation}


\int_\Omega ( {D (w), \eta} ) d \alpha 


+\int_\Omega ( \sqrt G \beta N(w), \eta ) d \alpha=0.


\label{In_Ur}


\end{equation}


Для $\alpha \in I^{-}(w) $ из (\ref{beta_2}) следует


\begin{equation}


(\sqrt G \beta N(w), \eta )=0.


\label{bzb6}


\end{equation}


Для $\alpha \in I(w)$, поскольку $\eta \in M(w)$


(и, значит, $w+t \eta \in M$ для $t \in [0,t_\eta ]$), имеем 


$F(w)+t \eta_3=w_3+t \eta_3 \geq F(w+t \eta)$,


откуда получаем


\begin{equation}


t \eta_3 \geq F(w+t \eta)-F(w)=t (\eta, \nabla F(w) )+o(t).


\label{ner1}


\end{equation}


Разделив неравенство (\ref{ner1}) на $t>0$ и


переходя при фиксированном $\alpha \in I(w) $ к 


пределу при $t \to+0$, получаем неравенство 


$( {[\tau_1,\tau_2], \eta} )=\eta_3 -(\eta, \nabla F(w) ) \geq 0$.


Из этого неравенства с учетом (\ref{Opr_N}) вытекает


\begin{equation}


( \sqrt G \beta N(w), \eta ) \geq 0,\ \ \alpha \in I(w).


\label{bzb7}


\end{equation} 





Таким образом, из (\ref{In_Ur}), (\ref{bzb6}), (\ref{bzb7}) следует, что 


если $w$ --- решение задачи (\ref{Ur_Rav_Pr}), то выполнено неравенство


\begin{equation} 


\int_\Omega ( {D (w(\alpha )), \eta (\alpha )} ) d \alpha 


\leq 0 \ \ \forall \eta \in M(w) . 


\label{Var_ner}


\end{equation}


Под вариационной постановкой этой задачи будем понимать


следующее: найти функцию $w$ из множества $M$, удовлетворяющую


условию (\ref{Var_ner}).


Предыдущие рассуждения показали, что если функция $w$ является решением


поточечной задачи (\ref{Ur_Rav_Pr}), то она является решением и вариационной


задачи.


Обратно, пусть функция $w$ из $M$ удовлетворяет вариационному


условию (\ref{Var_ner}). Определим функцию $\beta$ соотношением


\begin{equation}


\beta (\alpha )=\frac {| D(w(\alpha )) |}


{| [\partial_1 w(\alpha),\partial_2 w(\alpha) ] |},


\ \ \alpha \in \ov\Omega . 


\label{bzb9}


\end{equation}





Докажем, что для функций $w$ и $\beta$ выполнены уравнение


(\ref{Ur_Rav_Pr}) и условия (\ref{beta_1}), (\ref{beta_2}).


Убедимся, что 


\begin{equation} 


D(w(\alpha ))=- |D(w(\alpha )) | N(w(\alpha ))


\ \ \text{для}\ \ \alpha \in I(w) . 


\label{bzb10}


\end{equation}


Допустим, что нашлась точка $\alpha^* \in I(w) $, в которой не выполнено 


условие (\ref{bzb10}). Определим вектор $g=(g_1, g_2, g_3)$ равенством


$$


g=\frac {D(w^*) / |D(w^*)|+N(w^*)}


{|D(w^*) / |D(w^*)|+N(w^*)|},\ \ \text{где}\ \ w^*=w(\alpha^*).


$$


Этот единичный вектор удовлетворяет условиям


\begin{equation} 


(g, D(w^*) / |D(w^*) | ) > 0,\ \ (g, N(w^*) ) > 0.


\label{bzb11}


\end{equation}


Действительно, для любых двух единичных векторов $a$, $b$ таких, что 


выполнено условие $a \neq - b$, верно неравенство


$(a,b)>-1$. Eсли $c=(a+b)/|a+b|$, то $(a,c)>0$, $(b,c)>0$, поскольку


$$


(a, c)=\frac {(a, a+b)} {|a+b|}=\frac {1+(a, b)} {|a+b|}=(b, c) > 0 . 


$$ 


Левые части неравенств (\ref{bzb11}) по предположению, 


сделанному выше, непрерывны, 


а значит, найдутся такие $\delta, \varepsilon_0 > 0$, что


\begin{align} 


(g, D(w(\alpha)) ) &> \delta,\ \ \alpha \in B_{\varepsilon_0} (\alpha ^*), 


\label{bzb12} \\


%


( g, N(w(\alpha )) )|\nu (w(\alpha ))| &>


\delta,\ \ \alpha \in B_{\varepsilon_0} (\alpha ^*),


\label{bzb13}


\end{align}


где


$B_{\varepsilon}(\alpha)=\{\tau\in R^2: |\tau-\alpha|\leq\varepsilon \}$.


Далее, для ${\varepsilon > 0}$ введем функцию 


$\omega_{\varepsilon} \in C^{\infty} (R^2)$


со свойствами (см., напр., \cite{Vlad_1988}, с.\,89)


$$


0 \leq \omega _{\varepsilon} (\alpha ) \leq 1,


\ \ \omega _{\varepsilon} (\alpha )=1, 


\ \ \alpha \in B_{\varepsilon / 3}(0),


\ \ \omega _{\varepsilon} (\alpha )=0, 


\ \ \alpha \notin B_{\varepsilon}(0),


$$


и определим вектор-функцию 


$\eta (\alpha )=\varepsilon \omega _{\varepsilon} (\alpha -\alpha ^*) g$, 


$\alpha \in \Omega $.


Установим, что $\eta \in M(w)$ для достаточно малых $\varepsilon > 0$. 


При любых $t \in [0, 1]$, используя разложение в ряд Тейлора


функции $F$, имеем


\begin{multline}


w_3 (\alpha )+t \eta_3(\alpha ) - F(w (\alpha )+t \eta(\alpha ))=


w_3 (\alpha )+t \varepsilon \omega _{\varepsilon}(\alpha -\alpha ^*) g_3 - \\


%


-(F(w (\alpha ))+t \varepsilon \omega _{\varepsilon} (\alpha -\alpha ^*) 


( g, \nabla F(w( \alpha )))+


c(\alpha ) ( t \varepsilon \omega _{\varepsilon} (\alpha -\alpha ^*) )^2)=\\


%


=w_3 (\alpha ) - F(w (\alpha ))+


t \varepsilon \omega _{\varepsilon} (\alpha -\alpha ^*)


( ( N(w (\alpha )), g ) |\nu (w(\alpha ))|


- c(\alpha ) t \varepsilon 


\omega _{\varepsilon} (\alpha -\alpha ^*)),


\label{bzb14}


\end{multline}


где $c(\alpha )=(\nabla^2 F(x(\alpha )) g, g )$,


$x(\alpha)=w (\alpha )+\theta (\alpha ) \eta (\alpha )$,


$\theta (\alpha ) \in [0,1]$.


Предполагая, что $|c(\alpha )| \leq c_0$ при 


$\alpha \in B_{\varepsilon} (\alpha ^*)$, из (\ref{bzb13})


для $\varepsilon \leq \min \{\varepsilon _0$, $\delta / (2 c_0) \}$


получаем


$$


( N(w (\alpha )), g ) |\nu (w(\alpha ))|


- c(\alpha ) t \varepsilon 


\omega _{\varepsilon} (\alpha -\alpha ^*) \geq 


\delta - c_0 \delta / (2 c_0)=\delta / 2 > 0.


$$


Для $\alpha \notin B_{\varepsilon} (\alpha ^*)$


имеем $\omega _{\varepsilon} (\alpha -\alpha ^*)=0$, а значит, 


$$


t \varepsilon \omega _{\varepsilon} (\alpha -\alpha ^*) 


( {( N(w (\alpha )), g ) |\nu (w(\alpha ))|-c(\alpha ) t \varepsilon 


\omega _{\varepsilon} (\alpha -\alpha ^*)} )=0.


$$


Таким образом, поскольку $w \in M$


(т.\,е. $w_3(\alpha ) \geq F(w(\alpha ))$),


то из (\ref{bzb14}) следует, что при $t \in [0,1]$


$$


w_3 (\alpha )+t \eta_3(\alpha ) - F(w (\alpha )+t \eta(\alpha )) \geq 0,


\quad \alpha \in \ov\Omega,


$$


и поэтому $\eta \in M(w)$.


Для построенной функции $\eta$, пользуясь (\ref{bzb12}) и свойствами


функции $\omega _{\varepsilon}$, имеем


\begin{multline}


\int_\Omega ( {D (w(\alpha )), \eta (\alpha )} ) d \alpha=


\int_{B_{\varepsilon} (\alpha ^*) \cap \Omega} 


\varepsilon \omega _{\varepsilon} (\alpha -\alpha ^*)


( D (w(\alpha )), g ) d \alpha \geq \\


%


\geq \delta \varepsilon \int_{B_{\varepsilon / 3}(\alpha ^*) \cap \Omega} 


\omega _{\varepsilon} (\alpha -\alpha ^*) d\alpha=


\delta \varepsilon \mes (B_{\varepsilon / 3} (\alpha ^*) \bigcap \Omega )>0,


\label{bzb14_1}


\end{multline}


что противоречит (\ref{Var_ner}). Справедливость (\ref{bzb10}) тем


самым установлена.





Далее докажем равенство


\begin{equation} 


D(w(\alpha ))=0 \ \ \text{для}\ \ \alpha \in I^{-}(w) . 


\label{bzb10_1}


\end{equation}


Предположим, что нашлась точка


$\alpha^* \in I^{-}(w) $, в которой не выполнено 


условие (\ref{bzb10_1}). Положим $ g=D(w^*) / |D(w^*)|$, тогда 


найдутся такие $\delta, \varepsilon_0 > 0$, что выполнено (\ref{bzb12}).


Поскольку $\alpha^* \in I^{-}(w)$, то $w_3(\alpha^* )>F(w(\alpha^* )) $, и


в силу непрерывности функций $w$, $F$, найдется окрестность 


точки $\alpha^*$, в которой


$$


w_3 (\alpha ) - F(w (\alpha )) \geq \delta_1> 0.


$$


Пользуясь последним неравенством и (\ref{bzb14}),


получаем, что функция $\eta=\varepsilon \omega _{\varepsilon}g$


принадлежит $M(w)$, и для нее выполнено неравенство (\ref{bzb14_1}).


Получено противоречие с (\ref{Var_ner}) и тем


самым установлено равенство (\ref{bzb10_1}).





Из равенств (\ref{bzb9}), (\ref{bzb10})


и (\ref{bzb10_1}) следует, что функции $w$ и $\beta$


удовлетворяют уравнению (\ref{Ur_Rav_Pr})


и условиям (\ref{beta_1}), (\ref{beta_2}).


Таким образом, поточечная и вариационная


формулировки задачи (\ref{Ur_Rav_Pr}),


(\ref{Var_ner}) эквивалентны.





В рассмотренных задачах неизвестная функция $w$ описывает равновесное


положение оболочки. Сформулируем задачу относительно перемещений $u$


оболочки из известного первоначального положения $\xi$. Введем множества


\begin{gather} 


K=\{v : \ov\Omega \to R^3,


\ {v} |_\Gamma =0 ,\ \xi_3+v_3 \geq F(\xi+v), \alpha \in \ov\Omega \},


\label{K} \\


%


K(u)=\{{v \in K: \forall t \in [0,1],\ u+t ( v - u ) \in K} \} .


\label{Ku}


\end{gather}


Тогда вариационная задача относительно перемещений оболочки сводится 


к следующей: найти функцию $u$ из множества $K$, удовлетворяющую 


условию


\begin{equation} 


\int_\Omega ({D (\xi(\alpha )+u(\alpha)), v(\alpha )- u(\alpha )})d \alpha 


\leq 0 \ \ \forall v \in K(u) . 


\label{Var_ner_u}


\end{equation}


Докажем, что задачи (\ref{Var_ner}) и (\ref{Var_ner_u}) эквивалентны.


Действительно, пусть $u$ удовлетворяет задаче (\ref{Var_ner_u}), тогда


равновесное положение оболочки определено функцией $w=\xi+u$.


Выберем произвольную $\eta$ из $M(w)$, тогда $v=u+t_\eta\eta$


принадлежит множеству $K(u)$, и из (\ref{Var_ner_u}) получаем, что $w$


является решением задачи (\ref{Var_ner}). Аналогично доказывается, что


если $w$ --- решение вариационной задачи (\ref{Var_ner}), то 


$u=w-\xi$ --- решение (\ref{Var_ner_u}).





Заметим, что если множество $K$ выпуклое, то $K(u)=K$ для любого


$u$ из $K$, и квазивариационное неравенство (\ref{Var_ner_u})


становится вариационным.


Множество $K$ является выпуклым, если функция $\varphi$, описывающая


поверхность препятствия, выпукла во введенной декартовой системе координат. 





\section{Существование решения обобщенной задачи для сетчатой оболочки}





Рассмотрим обобщенную задачу в перемещениях об определении положения 


равновесия мягкой сетчатой оболочки, закрепленной по краям, 


находящейся под воздействием массовых 


сил, ограниченной препятствием,


и установим существование решения этой задачи.





Под сетчатой понимается оболочка, 


силовой основой которой является сетка, образованная двумя системами взаимно 


перекрещивающихся, абсолютно гибких упругих нитей. Предполагается,


что узлы сети фиксированы, материал, заполняющий промежутки 


между нитями, не сопротивляется деформации, и ни в начальном 


состоянии, ни в процессе деформации соседние нити не соприкасаются.


Ячейки сети считаются малыми и не сопротивляющимися сдвиговым деформациям.


Деформации и перемещения допускаются конечными.





Лагранжевы координаты $(\alpha^1,\alpha^2)$ выберем так, что координатные 


линии станут сонаправлены с нитями, образующими оболочку;


функция $\xi$ по-прежнему удовлетворяет условиям (\ref{2}).





Введем также следующие обозначения (для $k=1,2$):


$k^*=3-k$; $g_k=| \partial_k \xi |$, $G_{k}=| \partial_k w |$ ---


параметры Ламе поверхности 


недеформированной и деформированной оболочки соответственно; 


$R^1,R^2$ --- векторы, образующие контравариантный


локальный базис на деформированной поверхности:


$(R_k,R^m)=\delta_{km}$.





Обозначим через $\mathrm F^k$ внутреннюю силу, действующую


на единицу длины $\alpha^{k^*}$-й


координатной линии ($\alpha^k=\const$) 


деформированной оболочки,


с той стороны оболочки, куда направлен вектор $R^k$ ($k=1,2$), через


$F^{km}$ --- коэффициенты разложения этой плотности сил по единичным 


векторам локального базиса


$\mathrm F^k=\sum\limits_{m=1}^2 F^{km} R_m / G_m$.


Тогда компоненты тензора напряжений связаны с погонными усилиями 


$F^{km}$ соотношениями \cite{R_Gul}


$\sqrt G T^{km}=F^{km}G_{k^*}/G_m$.





Для сетчатой оболочки в силу того, что в выбранной лагранжевой системе


координат направления осей совпадают с направлениями нитей, имеем


(см. \cite{R_Gul}, \cite{B_Bux}) $F^{12}=F^{21}=0$


(т.\,е. ячейка сети не оказывает сопротивления повороту нитей 


в узлах скрепления),


$F^{kk}=N_k(\lambda_k)\rho_k g_{k^*}/G_{k^*}$,


где $\lambda_k=G_k/g_k$ --- относительные степени удлинения.





Здесь $N_1,N_2:R_+\to R_+$ --- функции, характеризующие физические свойства


нитей,\linebreak


$\rho_k:\Omega \to R_+$ --- количество нитей, сонаправленных с 


$\alpha^k$-й координатной осью, на единицу длины \linebreak


$\alpha^{k^*}$-й координатной оси в недеформированном состоянии. 


Эти функции определены конструкцией сетчатой 


оболочки, и относительно них считаем выполненными условия


\begin{gather}


N_k \in C_0(R_+),


\label{N1} \\


%


N_k(t)=0 \ \ \text{при}\ \ t \le 1 


\label{N2}


\end{gather}


(т.\,е. нити не воспринимают сжимающих усилий), 


\begin{equation}


N_k(t)>N_k(s)\ \ \text{при}\ \ t>s \ge 1, 


\label{N11}


\end{equation}


существуют $c,c_0,c_1,C > 0$, $p>1$ такие, что при $t \ge 0$


\begin{gather}


c_0 t^p - c_1 \le N_k (t)t \le C t^p, \label{N3} \\


%


\rho_k \in C_0 (\ov \Omega), \label{ro1} \\


%


\rho_k(\alpha) \geq c > 0 \ \ \forall\alpha\in \ov \Omega . \label{ro2}


\end{gather}





Заметим, что, вообще говоря, направления $R_k$ не являются


главными для тензора $\mathrm T$, хотя смешанные компоненты


$T^{km}$ ($k\neq m$) и равны нулю.





Плотность массовых сил $Q:\Omega \to R^3$ считаем известной. 


В силу закона сохранения массы имеем


$\sqrt G \gamma={| [\partial_1 \xi (\alpha),\partial_2 \xi (\alpha) ] |}


\overset{\circ}{\gamma}$,


где $\overset{\circ}{\gamma}{}:\Omega \to R^1$ ---


заданная плотность материала 


недеформированной оболочки. Относительно $Q$ и $\overset{\circ}{\gamma}$


считаем выполненными условия


$Q \in [C(\ov \Omega) ]^3$,


$\overset{\circ}{\gamma}{}\in C(\ov \Omega)$,


$\overset{\circ}{\gamma}{}(\alpha) \ge c >0$ 


\, $\forall \alpha\in \ov \Omega$.


Поверхностная нагрузка предполагается равной нулю: $P=0$.





Таким образом, в рассматриваемом случае


$$


%% \begin{equation}


D(w)=\sum_{k=1}^{2} \partial_k \bigg( \frac {N_k (|\partial_k w | / g_k )} 


{|\partial_k w |} \rho_k g_{k^*} \partial_k w \bigg)+


| [\partial_1 \xi,\partial_2 \xi ] | \overset{\circ}{\gamma}{}Q. 


%%\label{D_sell}


%%\end{equation}


$$





Введем пространство $V=[\overset{\circ}{W^1_p}{}(\Omega ) ]^3$


(где $p>1$ --- параметр из условия (\ref{N3})),


обозначим через $\langle \cdot,\cdot \rangle$ --- отношение двойственности 


между $V^*$ и $V$, где $V^*$ --- пространство, сопряженное к $V$.





Уточним для рассматриваемого случая сетчатой оболочки


определение множеств (\ref{K}), (\ref{Ku})


\begin{gather*}


K=\{v \in V: \xi_3(\alpha)+v_3(\alpha ) \geq F(\xi(\alpha)+v(\alpha )) 


\ \text{почти всюду на} \ \Omega \}, \\


%


K(u)=\{v \in K : \forall t \in [0,1],\ u+t ( v -u ) \in K \} .


\end{gather*}





Под решением обобщенной задачи (\ref{Var_ner_u}) будем понимать 


такую функцию $u \in K$, что для любого $v \in K(u)$


справедливо интегральное неравенство


\begin{equation}


\sum_{k=1}^{2} \int_\Omega


\frac{N_k (|\partial_k (\xi+u ) | / g_k )} {|\partial_k (\xi+u ) |} 


\rho_k g_{k^*} ( \partial_k (\xi+u ), \partial_k (v-u)) d\alpha \geq 


\int_\Omega


(|[\partial_1 \xi,\partial_2 \xi]| \overset{\circ}{\gamma}{}Q, v-u) d\alpha.


\label{Int_ner_sell}


\end{equation}


Зададим операторы $A_k: V \to V^{*}$ и функционал $f \in V^*$ формами


\begin{gather*} 


\langle A_k u, v \rangle=\int_\Omega \frac


{N_k (|\partial_k (\xi+u ) | / g_k )} {|\partial_k (\xi+u ) |} 


\rho_k g_{k^*} ( \partial_k (\xi+u), \partial_k v) d\alpha,\ \ k=1,2, \\


%


\langle f, v \rangle=\int_\Omega 


(| [\partial_1 \xi,\partial_2 \xi ] | \overset{\circ}{\gamma}{}Q, v) d\alpha. 


\end{gather*}





Убедимся в корректности определений. 


Из условий (\ref{2}) следует


\begin{equation}


C \ge g_k(\alpha) \geq c > 0 \ \ \forall\alpha\in \ov \Omega . 


\label{a2}


\end{equation}


Из (\ref{N3}), учитывая (\ref{ro1}), (\ref{a2}), получаем оценку


$$


N_k (|\partial_k u | / g_k ) \rho_k g_{k^*}


\le C |\partial_k u |^{p-1},\ \ \alpha\in \ov \Omega . 


$$


Из нее получаем корректность определения оператора $A_k$


\begin{multline*}


|\langle A_k u, v \rangle | \le C


\int_\Omega |\partial_k (u+\xi ) |^{p-1} |\partial_k v | d\alpha \le \\


%


\le C {\|\partial_k (u+\xi ) \|}_{L_p}^{p-1}{\|\partial_k v \|}_{L_p} 


\le C ({\| u \|}_V+{\| \partial_k \xi \|}_{L_p})^{p-1} {\| v \|}_V. 


%%\label{ogr_a} 


\end{multline*}





Пусть $A=A_1+A_2$. С учетом введенных обозначений обобщенная


задача (\ref{Int_ner_sell}) эквивалентна 


квазивариационному неравенству: 


найти $u$ из множества $K$, удовлетворяющую условию


\begin{equation}


\langle {A(u), v -u} \rangle \geq \langle {f, v-u} \rangle


\ \ \forall v \in K(u).


\label{Var_ner_sell} 


\end{equation}





\begin{thmnonum}


Задача $(\ref{Var_ner_sell})$ имеет решение.


\end{thmnonum}





\begin{pf}


Для рефлексивного банахова пространства $V$ 


по аналогии с \cite{BZ_92IV} из условий


(\ref{N1})--(\ref{ro2}) получаем, что операторы $A_k$, $k=1,2$, 


являются монотонными, деминепрерывными и коэрцитивными, а по


аналогии с \cite{LBK_78IV} устанавливается потенциальность оператора $A_k$.


Его потенциалом является функционал $\Phi_k: V \to R^1$


$$


\Phi_k(v)=\int_\Omega I_k (|\partial_k (\xi+v ) | / g_k) 


\rho_k g_{k^*} g_k d\alpha,\ \ \text{где}\ \ I_k(t)=\int_0^t N_k(t) dt .


$$


Определим функционал


$\Phi (v)=\Phi_1(v)+\Phi_2(v)-\langle f, v \rangle$, $v \in V$. 


В силу вышесказанного $\Phi $ является


слабо полунепрерывным снизу, коэрцитивным функционалом,


и ${\Phi}'(v)=Av - f$. 





Докажем теперь слабую замкнутость множества $K$. Пусть последовательность 


$\{v^{(i)} \}_{i=0}^{\infty}$, принадлежащая множеству $K$, слабо сходится


к $v$ в пространстве $V$.


Из компактности вложения пространства ${W^1_p}(\Omega )$ в $L_1(\Omega)$


следует, что последовательность $\{v^{(i)} \}_{i=0}^{\infty}$


сходится сильно к $v$ в пространстве $[L_1(\Omega ) ]^3$, а значит, найдется


подпоследовательность, сходящаяся почти всюду к $v$.


Поэтому $v$ принадлежит $K$.





Из установленных выше свойств функционала $\Phi$, 


множества $K$ и пространства $V$ 


следует по теореме Вейерштрасса \cite{Van72} существование решения задачи


минимизации функционала $\Phi$ на множестве $K$.


Пусть $u$ --- решение этой задачи


$$


\Phi (u)=\min_{v \in К} \Phi (v) .


$$


Докажем, что $u$ является решением квазивариационного 


неравенства (\ref{Var_ner_sell}). 


Пусть $v$ --- произвольный элемент множества $K(u)$. 


Из определения множества $K(u)$ следует, что 


$u+t(v-u) \in K$ для всех $t \in [0,1]$.


Элемент $u$ является точкой минимума $\Phi$ на множестве $K$, значит, 


\begin{equation}


\Phi (u+t(v-u) ) \ge \Phi (u),\ \ t \in [0,1].


\label{fun1} 


\end{equation}


В силу дифференцируемости $\Phi$ имеем


$$


\Phi(u+t(v-u))-\Phi (u)=t\langle{{\Phi}'(u+\theta_t t (v-u) ), v-u}\rangle,


\ \ \text{где}\ \ 0 \le \theta _t \le 1,


$$


следовательно, из (\ref{fun1}) получаем


$$


\langle {A(u+\theta _t t (v-u)) - f, v-u} \rangle \ge 0.


$$


Устремляя $t$ к нулю и учитывая деминепрерывность оператора $A$,


получаем неравенство (\ref{Var_ner_sell}).


\end{pf}





\begin{thebibliography}{99}





\bibitem{BS_86}


Бадриев И.Б., Шагидуллин Р.Р.


{\em Классические и обобщенные решения уравнений 


одноосного статического состояния мягкой оболочки} // Сеточ.


методы решения дифференц. уравнений. -- Казань: Изд-во Казанск. ун-та, 


1986. -- С.\,14--28. 





\bibitem{BZB_2001_2}


Бадриев И.Б., Задворнов О.А., Бандеров В.В.


{\em Постановка и исследование стационарных задач теории мягких оболочек


с невыпуклым достижимым множеством} // Исслед. по прикл.


матем. и информатике. -- Казань: Изд-во Казанск. матем. о-ва, 2001. --


Вып.\,23. -- С.\,3--7.





\bibitem{BBZ_2002_1}


Бадриев И.Б., Бандеров В.В., Задворнов О.А.


{\em Задача о равновесном положении нити с невыпуклым 


достижимым множеством} // Актуальн. проблемы матем. моделирования


и информатики. -- Казань: Изд-во Казанск. матем. о-ва, 2002. --


С.\,25--29.





\bibitem{BZ_92IV}


Бадриев И.Б., Задворнов О.А.


{\em Исследование разрешимости стационарных задач для 


сетчатых оболочек} // Изв. вузов. Математика.


-- 1992. -- \No\,11. -- С.\,3--7.





\bibitem{R_Gul}


Ридель В.В., Гулин Б.В.


{\em Динамика мягких оболочек}. -- М.: Наука, 1990. -- 206~с.





\bibitem{Duv}


Дюво Г., Лионс Ж.-Л.


{\em Неравенства в механике и физике}. --


М.: Наука, 1980. -- 383 с.





\bibitem{Pan}


Панагиотопулос П.


{\em Неравенства в механике и их приложения. Выпуклые и невыпуклые 


функции энергии}. -- М.: Мир, 1989. -- 496 с.





\bibitem{Kid}


Kydoniefs A.D.


{\em Finite axisymmetric deformations of an initially cylindrical 


elastic membrane enclosing a rigid body} //


Quart. Journ. Mech. and Appl. Math. -- 1969. -- V.\,XXII.


-- P.\,319--331.





\bibitem{Vlad_1988}


Владимиров В.С. {\em Уравнения математической физики}. -- М.: Наука, 


1988. -- 512~ с. 





\bibitem{B_Bux}


Бидерман В.Л., Бухин Б.Л.


{\em Уравнения равновесия безмоментной сетчатой оболочки} //


Инж. журн. МТТ. -- 1966. -- \No\,1. -- С.\,84--89. 





\bibitem{LBK_78IV}


Ляшко А.Д., Бадриев И.Б., Карчевский М.М.


{\em О вариационном методе для уравнений с разрывными 


монотонными операторами} // Изв. вузов. Математика.


-- 1978. -- \No\,11. -- С.\,63--69. 





\bibitem{Van72}


Вайнберг М.М. {\em Вариационный метод и метод монотонных операторов}. --


М.: Наука, 1972. -- 415 с.





\end{thebibliography}





\vskip 2cm





\post{\it Казанский государственный}{\it Поступила}


\post{\it университет}{20.06.2002}





\label{end}


\end{document}


